\subsection{\texorpdfstring{THE COEFFICIENTS \(c_{\alpha\beta\gamma}\)}{THE COEFFICIENTS c\_\{\textbackslash alpha\textbackslash beta\textbackslash gamma\}}}

Let \(\Omega\) be the set of \((x,y)\in\mathbf{U}\times\mathbf{U}\) such that \(x.y\) is defined and belongs to \(\mathbf{U}\) . Then \(\Omega\) is open in \(\mathbf{U}\times\mathbf{U}\) and the mapping \((x,y)\mapsto x.y\) of \(\Omega\) into \(\mathbf{U}\) is analytic. The coordinates \(z_{1},\ldots,z_{n}\) of \(z=x.y\) therefore admit expansions as integral series about the origin in the powers of \(x_{1},\ldots,x_{n},y_{1},\ldots,y_{n}\) . Therefore, there exist well defined constants \(c_{\alpha_{1},\ldots,\alpha_{n},\beta_{1},\ldots,\beta_{n}}\in\mathbf{K}\) , such that

\[
z _ {1} ^ {\gamma_ {1}} \dots z _ {n} ^ {\gamma_ {n}} = \sum_ {\alpha_ {1},\dots,\alpha_ {n},\beta_ {1},\dots,\beta_ {n} \in \mathbf {N}} c _ {\alpha_ {1},\dots,\alpha_ {n},\beta_ {1},\dots,\beta_ {n},\gamma_ {1},\dots,\gamma_ {n}} x _ {1} ^ {\alpha_ {1}} \dots x _ {n} ^ {\alpha_ {n}} y _ {1} ^ {\beta_ {1}} \dots y _ {n} ^ {\beta_ {n}}\tag{1}
\]

for \(\gamma_{1},\ldots,\gamma_{n}\) in \(\mathbf{N}\). Adopting the conventions of Differentiable and Analytic Manifolds, R, we shall write these formulae more briefly:

\[
(x. y) ^ {\gamma} = \sum_ {\alpha , \beta \in \mathbf {N} ^ {n}} c _ {\alpha , \beta , \gamma} x ^ {\alpha} y ^ {\beta} \quad (\gamma \in \mathbf {N} ^ {n}).\tag{2}
\]

Since \(x.0 = 0.x = x\) for \(x\in \mathbf{U}\)

\[
c _ {\alpha , 0, \gamma} = c _ {0, \alpha , \gamma} = \delta_ {\alpha \gamma}\tag{3}
\]

where \(\delta_{\alpha \gamma}\) is the Kronecker index. In particular, writing henceforth \(k\) instead of \(\varepsilon_k\) for \(k = 1, \ldots, n\) ,

\[
(x. y) _ {k} = x _ {k} + y _ {k} + \sum_ {| \alpha | \geqslant 1, | \beta | \geqslant 1} c _ {\alpha , \beta , k} x ^ {\alpha} y ^ {\beta}.\tag{4}
\]

Writing \(c_{\alpha \beta} = (c_{\alpha \beta 1}, c_{\alpha \beta 2}, \ldots, c_{\alpha \beta n}) \in \mathbf{K}^n\) , it then follows that

\[
x. y = x + y + \sum_ {| \alpha | \geqslant 1, | \beta | \geqslant 1} c _ {\alpha \beta} x ^ {\alpha} y ^ {\beta}.\tag{5}
\]

On the right hand side of (5), we consider the homogeneous component of degree 2:

\[
\mathrm{B} (x, y) = \sum_ {i, j = 1} ^ {n} c _ {i j} x _ {i} y _ {j}
\]

so that \((x,y)\mapsto \mathrm{B}(x,y)\) is a bilinear mapping of \(\mathbf{K}^n\times \mathbf{K}^n\) into \(\mathbf{K}^n\) . Then

\[
x. y \equiv x + y + \mathrm{B} (x, y) \mod \deg 3.\tag{6}
\]

Formula (4) implies on the other hand that

\[
c _ {\alpha , \beta , \gamma} = 0 \quad \text { if } | \alpha | + | \beta | <   | \gamma |\tag{7}
\]

and that the terms of total degree \(|\gamma|\) in the expansion of \(z^\gamma\) are also those of \((x_1 + y_1)^{\gamma_1}(x_2 + y_2)^{\gamma_2} \ldots (x_n + y_n)^{\gamma_n} = \sum_{\alpha + \beta = \gamma} ((\alpha, \beta)) x^\alpha y^\beta\) (cf.~Differentiable and Analytic Manifolds, R, Notation and conventions). Hence:

\begin{enumerate}
\def\labelenumi{(\arabic{enumi})}
\setcounter{enumi}{7}
\tightlist
\item
\end{enumerate}

\[
c _ {\alpha , \beta , \gamma} = 0 \quad \text { if } \quad | \alpha | + | \beta | = | \gamma | \quad \text { but } \quad \alpha + \beta \neq \gamma\tag{8}
\]

\[
c _ {\alpha , \beta , \alpha + \beta} = ((\alpha , \beta)).\tag{9}
\]

The associativity of the product implies the relation

\[
\sum_ {\alpha , \xi} c _ {\alpha \xi \eta} x ^ {\alpha} \left(\sum_ {\beta , \gamma} c _ {\beta \gamma \xi} y ^ {\beta} z ^ {\gamma}\right) = \sum_ {\xi , \gamma} c _ {\xi \gamma \eta} \left(\sum_ {\alpha , \beta} c _ {\alpha \beta \xi} x ^ {\alpha} y ^ {\beta}\right) z ^ {\gamma}
\]

for all \(x, y, z\) sufficiently close to 0, whence

\[
\sum_ {\xi} c _ {\alpha \xi \eta} c _ {\beta \gamma \xi} = \sum_ {\xi} c _ {\xi \gamma \eta} c _ {\alpha \beta \xi} (\alpha , \beta , \gamma , \eta \text { in } \mathbf {N} ^ {n}).\tag{10}
\]

The group germ G admits a commutative open subgroup germ if and only if \(c_{\alpha\beta\gamma} = c_{\beta\alpha\gamma}\) for all \(\alpha, \beta, \gamma\) in \(\mathbf{N}^{n}\) .

